\documentclass[12pt]{article}
\usepackage[a4paper,margin=1in]{geometry}
\usepackage{graphicx}
\usepackage{hyperref}
\usepackage{amsmath}

\newcommand{\Mod}[1]{\ (\mathrm{mod}\ #1)}

\title{Integer Modulo on Bitvectors is Weird}
\author{Dominic Kennedy}
\date{\today}

\begin{document}

\maketitle

\section{Introduction}
There is an obvious construction for the lattices of the known bits and constant range domains which is based on the size of the concrete bitvectors being abstracted over.
The known bits lattice is a ternary bitvector with the same width as the concrete bitvectors.
E.g. if the concrete bitvectors are 4 bits wide, then there are $2^4 = 16$ possible concrete values, and the known bits lattice has $3^{bw} = 3^4 = 81$ elements.
Similarly, a constant range domains lattice will always have $\frac{2^{bw}(2^{bw}+1)}{2}$ lattice values.

The integer modulo domain is a congruence relation between the bitvector and the first $n$-primes.
Unlike known bits and constant range,
in integer modulo $n$ may be selected to be as small as 0, or large enough such that the $n^{th}$-prime in the set is still representable by the bitvector
(i.e. smaller than $2^{bw}$).

\section{Main Content}

This selection of $n$ affects whether or not the lattice is well behaved.
With $2^{bw}$ being the distinct number of concrete values, and $P$ being the product of the first $n$-primes.
Either, $2^{bw} < P$, or $2^{bw} \geq P$.

\subsection{If $2^{bw} < P$}

There are two types of bad lattice values that arise when $2^{bw} < P$.
I have no proof that these are the only problems that show up, but after testing with various values for $bw$ and $P$ only these problems show up.

\subsubsection{Lattice Values Representing $\bot$}

Consider an integer modulo domain on the first three primes (2, 3, and 5) so $P = 30$ and $bw = 4$ so there are 16 distinct concrete values.
I'll show abstract values like this: $[1, 2, 1]$ which represents any integer $x$ such that
$x < 2^{bw}$, $x \equiv 1 \mod 2$, $x \equiv 2 \mod 3$, $x \equiv 1 \mod 5$, in this case $\gamma([1, 2, 1])=\{11\}$.
Abstract values written like this: $[1, \top, 1]$, mean that there is no congruence relation on mod 3, so in this case $\gamma([1, \top, 1])=\{1, 11\}$.

Now, consider the abstract value $[1, 1, 0]$, the first integer to satisfy the congruence relations is 25, but that's larger than our bitvectors represent.
So this abstract value represens $\bot$, in fact there are exactly $P - 2^{bw} = 14$ ways to representation $\bot$.
This is a problem since we want all lattice representations to be unique.

\subsubsection{Multiple Lattice Representations for singleton values}
Consider, $\gamma([0, 1, 4]) = \{4\}$.
But, $\gamma([\top, 1, 4]) = \{4\}$ as well, which once again means that there are multiple non-unique lattice values.
These duplicate values only occur when the concrete set is a single value (again I have no proof of this, other than testing with large $P$).

\subsubsection{Solutions}
I belive these to be the only two types of duplicate lattice values, that occur.
There's likely some cool number theory tom-foolery which would prove this, but I just checked with a wack ton of primes.

It's easy(-ish) to ``skip'' these problamatic lattice values in the \texttt{enumVals} function,
however its common for these ``bad'' lattice values to show up as the output of a transfer function.
Any transfer function which used these ``bad'' lattice values could still be technically 100\% precise,
but would certainly not be scored that way by the eval engine.
Further, if a transfer function produced these lattice values it could lead to lower precision when combined with other transfer functions.

The best solution, barring some number theory voodoo, is to call $\gamma$ every time an abstract value is created to test if it's in a ``bad'' state.
Then fix it if it does happen to be in such a state.
The problem with this solution is that it would make the eval engine (and any compiler which would use our synthesised transfer functions) too slow.

\subsection{If $2^{bw} \geq P$}

If big $P$ makes the lattice unreliable, then just reduce $P$.
As far as I know there are no problems with the lattice misbahving in this case.
It's just that the smaller lattice reduces abstract interpretation precision.

\subsubsection{some notes on lattice size vs powerset size}
As the width of a bitvector increases the number of concrete values it can represent grows at $\mathcal{O}(2^n)$.
The complete lattice of the bitvector grows at $\mathcal{O}(2^{2^n})$,
The known bits lattice grows at $\mathcal{O}(3^n)$.
the constant range lattice grows at $\mathcal{O}(4^n)$,
but integer modulo must be constructed such that $2^{bw} \geq P$
and the size of the lattice is only $\prod_{\forall p \in P} p+1$,
which means that it only grows at $\mathcal{O}(2^n)$ (the same rate as the number of possible concrete values themselves!!).
This makes it worse for precise abstract representations (I think).

\end{document}
